% ===== Джерела до розділу 12. Ключі з префіксом ch12: =====
% Позиції ядра (bib-core.tex) тут НЕ дублюються.

% --- Розвідка загроз: концепції та стандарти ---
\bibitem{ch12:mcmillan} McMillan R. Definition: Threat Intelligence. Gartner Research, 2013.

\bibitem{ch12:stix} OASIS. STIX Version 2.1: Structured Threat Information Expression. OASIS Standard, 2021. URL: \url{https://docs.oasis-open.org/cti/stix/v2.1/stix-v2.1.html}.

\bibitem{ch12:taxii} OASIS. TAXII Version 2.1: Trusted Automated Exchange of Intelligence Information. OASIS Standard, 2021.

\bibitem{ch12:misp} Wagner C., Dulaunoy A., Wagener G., Iklody A. MISP --- The Design and Implementation of a Collaborative Threat Intelligence Sharing Platform. \emph{Proceedings of the 2016 ACM on Workshop on Information Sharing and Collaborative Security (WISCS~'16)}. New York: ACM, 2016. P.~49--56.

% --- Моделі аналізу та атрибуція ---
\bibitem{ch12:diamond} Caltagirone S., Pendergast A., Betz C. The Diamond Model of Intrusion Analysis. Center for Cyber Intelligence Analysis and Threat Research, 2013.

\bibitem{ch12:rid} Rid T., Buchanan B. Attributing Cyber Attacks. \emph{Journal of Strategic Studies}. 2015. Vol.~38, №~1--2. P.~4--37. DOI: 10.1080/01402390.2014.977382.

\bibitem{ch12:attribution} Bartholomew B., Guerrero-Saade J. A. Wave Your False Flags! Deception Tactics Muddying Attribution in Targeted Attacks. \emph{Virus Bulletin Conference}. Denver, 2016.

% --- Кампанії та кейси ---
\bibitem{ch12:wannacry} Nakashima E., Timberg C. NSA Officials Worried about the Day Its Potent Hacking Tool Would Get Loose. Then It Did. \emph{The Washington Post}. 2017. (аналіз кампанії WannaCry).

\bibitem{ch12:greenberg} Greenberg A. Sandworm: A New Era of Cyberwar and the Hunt for the Kremlin's Most Dangerous Hackers. New York: Doubleday, 2019. 348~p.

% --- Українські джерела ---
\bibitem{ch12:certua} Команда реагування на комп'ютерні надзвичайні події України CERT-UA. Аналітичні звіти та індикатори компрометації. Державна служба спеціального зв'язку та захисту інформації України. URL: \url{https://cert.gov.ua/}.

\bibitem{ch12:ti-ua} Застосування розвідки кіберзагроз (threat intelligence) для проактивного захисту інформаційних систем. \emph{Кібербезпека: освіта, наука, техніка}. 2024. URL: \url{https://csecurity.kubg.edu.ua/index.php/journal/article/view/759}.

\bibitem{ch12:apt-ua} Аналіз діяльності APT-угруповань, спрямованих проти критичної інфраструктури України. \emph{Захист інформації}. 2023. Т.~25, №~2.
